% Generated from lessons/0009-5-infinite-products.html; edit the HTML source, then rerun the converter.
\section{无穷乘积}
\lessonmark{无穷乘积}

这一节表面上是“连乘”，本质上仍然是极限问题。最重要的桥梁是：把乘积写成 \(\prod(1+a_n)\) 后取对数，就会回到你已经熟悉的级数世界。定义、必要条件、对数判别、Wallis 公式和 \(\sin x\) 的无穷乘积，是这一节的五个核心点。

\subsection*{本节主线}

\begin{conceptbox}
\textbf{先盯部分积}
无穷乘积是否“收敛”，只看部分积 \(P_n=\prod_{k=1}^n p_k\) 是否趋于非零有限数。
\end{conceptbox}

\begin{conceptbox}
\textbf{再过对数桥}
当 \(p_n>0\) 时，\(\prod p_n\) 的问题可转成 \(\sum \ln p_n\) 的问题。
\end{conceptbox}

\begin{conceptbox}
\textbf{最后分类讨论}
同号项直接比 \(\sum a_n\)，变号项要多看一步 \(\sum a_n^2\)，绝对收敛则回到 \(\sum |a_n|\)。
\end{conceptbox}

\begin{notebox}
考试顺序：先写 \(p_n=1+a_n\)，再检查 \(a_n\to0\)；若最终同号，直接看 \(\sum a_n\)；若 \(\sum a_n\) 已知收敛，再看 \(\sum a_n^2\)；涉及具体值时优先想 Wallis 公式或裂项相消。
\end{notebox}

\subsection*{一、无穷乘积的定义}

\subsubsection*{定义 9.5.1}

设 \(p_1,p_2,\ldots,p_n,\ldots\)（每个 \(p_n\neq0\)）是一列实数。把

\[
      p_1\cdot p_2\cdot \cdots \cdot p_n\cdots
    \]

称为无穷乘积，记作

\[
      \prod_{n=1}^{\infty}p_n.
    \]

对应的部分积数列为

\[
      P_1=p_1,\quad P_2=p_1p_2,\quad \ldots,\quad P_n=\prod_{k=1}^{n}p_k.
    \]

如果 \(\{P_n\}\) 收敛到一个非零有限数 \(P\)，就称无穷乘积收敛，且

\[
      \prod_{n=1}^{\infty}p_n=P.
    \]

如果 \(\{P_n\}\) 发散，或 \(\{P_n\}\) 收敛于 \(0\)，都称无穷乘积发散。

\begin{notebox}
这里和级数最大的区别是：乘积“趋于 0”不叫收敛，而叫“发散于 0”。因为我们关心的是能否稳定停在一个非零值上。
\end{notebox}

\subsection*{二、收敛的必要条件}

\subsubsection*{定理 9.5.1}

若无穷乘积 \(\prod_{n=1}^{\infty}p_n\) 收敛，则：

\[
      (1)\ \lim_{n\to\infty}p_n=1;
      \qquad
      (2)\ \lim_{m\to\infty}\prod_{n=m+1}^{\infty}p_n=1.
    \]

\subsubsection*{证明思路}

若 \(P_n\to P\neq0\)，则

\[
      p_n=\frac{P_n}{P_{n-1}}\to \frac{P}{P}=1.
    \]

尾积也同理：

\[
      \prod_{n=m+1}^{\infty}p_n=\frac{P}{P_m}\to 1.
    \]

常把 \(p_n\) 写成

\[
      p_n=1+a_n.
    \]

于是定理 9.5.1 的第一个结论就变成：

\[
      \prod_{n=1}^{\infty}(1+a_n)\text{ 收敛}\quad\Longrightarrow\quad a_n\to0.
    \]

\begin{conceptbox}
\textbf{例 9.5.1}
\(\displaystyle p_n=1-\frac{1}{n+1}\)。则

\[
          P_n=\prod_{k=1}^{n}\left(1-\frac{1}{k+1}\right)
          =\frac12\cdot\frac23\cdots\frac{n}{n+1}
          =\frac{1}{n+1}\to0.
        \]

所以这个无穷乘积发散于 \(0\)。
\end{conceptbox}

\begin{conceptbox}
\textbf{必要条件的用途}
若 \(p_n\not\to1\)，可以立刻判发散，不用再做任何复杂计算。
\end{conceptbox}

\subsection*{三、Wallis 公式与基础例子}

\subsubsection*{例 9.5.2：Wallis 公式}

设

\[
      p_n=1-\frac{1}{(2n)^2},\qquad n=1,2,\ldots
    \]

则无穷乘积

\[
      \prod_{n=1}^{\infty}\left(1-\frac{1}{(2n)^2}\right)
    \]

收敛，并且

\[
      \prod_{n=1}^{\infty}\left(1-\frac{1}{(2n)^2}\right)=\frac{2}{\pi}.
    \]

换个形式，就是著名的 Wallis 公式：

\[
      \frac{\pi}{2}
      =
      \frac21\cdot\frac23\cdot\frac43\cdot\frac45\cdot\frac65\cdot\frac67\cdots
      =
      \prod_{n=1}^{\infty}\frac{4n^2}{4n^2-1}.
    \]

\subsubsection*{证明入口}

令

\[
      I_n=\int_0^{\pi/2}\sin^n x\,dx.
    \]

利用递推可得

\[
      I_{2n}=\frac{(2n-1)!!}{(2n)!!}\cdot\frac{\pi}{2},
      \qquad
      I_{2n+1}=\frac{(2n)!!}{(2n+1)!!}.
    \]

再用

\[
      I_{2n+1}<I_{2n}<I_{2n-1}
    \]

和夹逼定理，就能得到 \(\frac{I_{2n}}{I_{2n+1}}\to1\)，最后整理出 Wallis 公式。

\subsubsection*{例 9.5.3：余弦乘积}

设 \(p_n=\cos \dfrac{x}{2^n}\)，则当 \(0<x<\pi\) 时，

\[
      \prod_{n=1}^{\infty}\cos\frac{x}{2^n}=\frac{\sin x}{x}.
    \]

令 \(x=\dfrac{\pi}{2}\)，便得到 Viète 公式：

\[
      \frac{2}{\pi}
      =
      \cos\frac{\pi}{4}\cdot\cos\frac{\pi}{8}\cdot\cos\frac{\pi}{16}\cdots.
    \]

\begin{notebox}
这一节里最常用的具体公式，就是 Wallis 和 \(\prod \cos \frac{x}{2^n}=\frac{\sin x}{x}\)。遇到带 \(\sin x\)、\(\pi\)、双阶乘的题，先想这两个。
\end{notebox}

\subsection*{四、无穷乘积与无穷级数}

由定理 9.5.1 可知，若无穷乘积 \(\prod p_n\) 收敛，则从某项起 \(p_n>0\)。因此讨论收敛性时，可以默认最终都能取对数。

\subsubsection*{定理 9.5.2}

无穷乘积

\[
      \prod_{n=1}^{\infty}p_n
    \]

收敛的充要条件是级数

\[
      \sum_{n=1}^{\infty}\ln p_n
    \]

收敛。

\subsubsection*{证明思路}

设 \(\sum\ln p_n\) 的部分和为 \(S_n\)。由于

\[
      P_n=\prod_{k=1}^{n}p_k=e^{S_n},
    \]

\(\{P_n\}\) 收敛到非零有限数当且仅当 \(\{S_n\}\) 收敛到有限数。这就是无穷乘积和无穷级数之间的桥。

\subsubsection*{推论 9.5.1：同号情形}

若 \(a_n>0\) 或 \(a_n<0\)（从某项起保持同号），则

\[
      \prod_{n=1}^{\infty}(1+a_n)\text{ 收敛}
      \quad\Longleftrightarrow\quad
      \sum_{n=1}^{\infty}a_n\text{ 收敛}.
    \]

理由是 \(a_n\to0\) 时，\(\ln(1+a_n)\sim a_n\)。

\subsubsection*{推论 9.5.2：变号情形}

若级数 \(\sum a_n\) 收敛，则无穷乘积

\[
      \prod_{n=1}^{\infty}(1+a_n)
    \]

收敛的充要条件是级数

\[
      \sum_{n=1}^{\infty}a_n^2
    \]

收敛。

\subsubsection*{为什么会出现 \(a_n^2\)}

当 \(\sum a_n\) 已经收敛时，\(\ln(1+a_n)\) 和 \(a_n\) 的差异主要落在二次项：

\[
      \ln(1+a_n)=a_n-\frac{a_n^2}{2}+o(a_n^2).
    \]

所以对 \(\sum \ln(1+a_n)\) 来说，真正决定成败的是 \(\sum a_n^2\)。

\subsection*{五、绝对收敛}

\subsubsection*{例 9.5.4}

讨论无穷乘积

\[
      \prod_{n=1}^{\infty}\left(1+\frac{(-1)^{n+1}}{n^x}\right)
    \]

的敛散性。

\begin{center}
\small
\begin{tabularx}{\linewidth}{|>{\raggedright\arraybackslash}p{0.18\linewidth}|>{\raggedright\arraybackslash}p{0.42\linewidth}|>{\raggedright\arraybackslash}X|}
\hline
\textbf{参数} & \textbf{结论} & \textbf{理由} \\ \hline
\(x\le0\) & 发散 & 必要条件 \(a_n\to0\) 不成立 \\ \hline
\(0<x\le\frac12\) & 发散 & \(\sum a_n\) 收敛，但 \(\sum a_n^2=\sum 1/n^{2x}\) 发散 \\ \hline
\(x>\frac12\) & 收敛 & \(\sum a_n\) 与 \(\sum a_n^2\) 都收敛 \\ \hline
\end{tabularx}
\end{center}

\subsubsection*{定义 9.5.2}

当级数 \(\sum |\ln p_n|\) 绝对收敛时，称无穷乘积 \(\prod p_n\) 绝对收敛。

\subsubsection*{定理 9.5.3}

设 \(a_n>-1\)。下述三个命题等价：

\begin{enumerate}
 \item 无穷乘积 \(\prod_{n=1}^{\infty}(1+a_n)\) 绝对收敛；
 \item 无穷乘积 \(\prod_{n=1}^{\infty}(1+|a_n|)\) 收敛；
 \item 级数 \(\sum_{n=1}^{\infty}|a_n|\) 收敛。
 \end{enumerate}

所以对考试来说，可以直接记成：

\[
      \prod(1+a_n)\text{ 绝对收敛}
      \quad\Longleftrightarrow\quad
      \sum |a_n|\text{ 收敛}.
    \]

\begin{notebox}
把例 9.5.4 往前再推一步：\(\prod(1+\frac{(-1)^{n+1}}{n^x})\) 在 \(x>1\) 时绝对收敛，在 \(\frac12<x\le1\) 时只是非绝对收敛。
\end{notebox}

\subsection*{六、Stirling 公式与两个重要例题}

\subsubsection*{例 9.5.5：Stirling 公式}

证明

\[
      n!\sim \sqrt{2\pi}\,n^{n+\frac12}e^{-n}\qquad(n\to\infty).
    \]

教材做法是设

\[
      b_n=\frac{n!\,e^n}{n^{n+\frac12}}.
    \]

然后考察

\[
      \frac{b_n}{b_{n-1}}
      =
      e\left(1-\frac1n\right)^{n-\frac12}
      =
      1-\frac{1}{12n^2}+o\!\left(\frac1{n^2}\right).
    \]

于是 \(1+a_n=\frac{b_n}{b_{n-1}}\)，并且 \(\sum a_n\) 收敛，所以相应无穷乘积收敛到非零常数。再利用 Wallis 公式，可算出该常数正是 \(\sqrt{2\pi}\)。

\subsubsection*{例 9.5.6}

求极限

\[
      \lim_{n\to\infty}\frac{n}{\sqrt[n]{n!}}.
    \]

由 Stirling 公式：

\[
      n!\sim \sqrt{2\pi}\,n^{n+\frac12}e^{-n},
    \]

取 \(n\) 次方根得到

\[
      \sqrt[n]{n!}\sim \sqrt[n]{\sqrt{2\pi}\,n^{n+\frac12}e^{-n}}\sim \frac{n}{e},
    \]

故

\[
      \lim_{n\to\infty}\frac{n}{\sqrt[n]{n!}}=e.
    \]

\subsubsection*{例 9.5.7：\(\sin x\) 的无穷乘积展开}

证明

\[
      \sin x=x\prod_{n=1}^{\infty}\left(1-\frac{x^2}{n^2\pi^2}\right).
    \]

这是本节最漂亮的公式之一。教材用三角恒等式、根的结构和无穷乘积收敛性逐步推出。你可以把它记成“\(\sin x\) 的零点在 \(\pm n\pi\)，所以展开式里自然出现 \(1-\frac{x^2}{n^2\pi^2}\)”。

\subsection*{七、判题路线图}

\begin{center}
\small
\begin{tabularx}{\linewidth}{|>{\raggedright\arraybackslash}p{0.18\linewidth}|>{\raggedright\arraybackslash}p{0.42\linewidth}|>{\raggedright\arraybackslash}X|}
\hline
\textbf{看到什么} & \textbf{先试什么} & \textbf{提醒} \\ \hline
\(\prod p_n\) & 先看 \(p_n\to1\) 吗 & 不趋于 \(1\) 直接发散。 \\ \hline
\(\prod(1+a_n)\) 且 \(a_n\) 最终同号 & 比较 \(\sum a_n\) & 这是最快的一类题。 \\ \hline
\(\sum a_n\) 已知收敛 & 再看 \(\sum a_n^2\) & 这是变号题的核心分界。 \\ \hline
要判绝对收敛 & 看 \(\sum |a_n|\) & 等价于 \(\sum |\ln(1+a_n)|\) 绝对收敛。 \\ \hline
题里出现双阶乘、\(\pi\)、\(\sin x\) & 想 Wallis 或 \(\sin x\) 乘积 & 通常不是纯比较题。 \\ \hline
\end{tabularx}
\end{center}

\subsection*{八、即时判断练习}

\begin{quickcheck}
\[
          \prod_{n=1}^{\infty}\left(1-\frac{1}{n+1}\right)
        \]

\tcblower
\textbf{答案：}zero
\end{quickcheck}

\begin{quickcheck}
\[
          \prod_{n=1}^{\infty}\left(1+\frac{1}{n^2}\right)
        \]

\tcblower
\textbf{答案：}conv
\end{quickcheck}

\begin{quickcheck}
\[
          \prod_{n=1}^{\infty}\left(1+\frac{(-1)^{n+1}}{n^x}\right)
          \text{ 的临界指数是？}
        \]

\tcblower
\textbf{答案：}half
\end{quickcheck}

\subsection*{九、课后题}

下面按教材第九章第 5 节习题整理。题面完整保留，提示只给入口；写作业时建议先独立判断，再展开提示核对方向。

\begin{exercisebox}
\textbf{习题 1 \badge{讨论无穷乘积的敛散性}}

讨论下列无穷乘积的敛散性：

\[
      \begin{aligned}
      &(1)\ \prod_{n=1}^{\infty}\frac{n^2}{n^2+1};\qquad
      (2)\ \prod_{n=2}^{\infty}\sqrt{\frac{n+1}{n-1}};\\
      &(3)\ \prod_{n=3}^{\infty}\cos\frac{\pi}{n};\qquad
      (4)\ \prod_{n=1}^{\infty}n\sin\frac1n;\\
      &(5)\ \prod_{n=1}^{\infty}e^{1/n^2};\qquad
      (6)\ \prod_{n=1}^{\infty}\left(1-\frac{x^2}{n^2\pi^2}\right);\\
      &(7)\ \prod_{n=1}^{\infty}\left(1+\frac{x^n}{2^n}\right);\qquad
      (8)\ \prod_{n=1}^{\infty}\sqrt{1+\frac1n};\\
      &(9)\ \prod_{n=1}^{\infty}\left[\left(1+\frac{x}{n}\right)e^{-x/n}\right];\qquad
      (10)\ \prod_{n=1}^{\infty}\left[\left(1+\frac1{n^p}\right)\cos\frac{\pi}{n^q}\right]\quad(p,q>0).
      \end{aligned}
      \]

\begin{hintbox}
先统一写成 \(\prod(1+a_n)\)。第 (1)(3)(4)(5)(8) 先看 \(\sum a_n\)；第 (6) 可用 \(\sin x\) 的无穷乘积；第 (9) 取对数后会出现 \(x/n-\ln(1+x/n)\sim x^2/(2n^2)\)。
\end{hintbox}
\end{exercisebox}

\begin{exercisebox}
\textbf{习题 2 \badge{计算无穷乘积的值}}

计算下列无穷乘积的值：

\[
      \begin{aligned}
      &(1)\ \prod_{n=2}^{\infty}\left(1-\frac1{n^2}\right);\\
      &(2)\ \prod_{n=2}^{\infty}\left(1-\frac{2}{n(n+1)}\right);\\
      &(3)\ \prod_{n=2}^{\infty}\frac{n^3-1}{n^3+1}.
      \end{aligned}
      \]

\begin{hintbox}
这三题都适合先写部分积。第 (1) 可直接裂项；第 (2)(3) 先因式分解，再观察相消结构。
\end{hintbox}
\end{exercisebox}

\begin{exercisebox}
\textbf{习题 3 \badge{余弦乘积与平方和}}

设 \(0<x_n<\dfrac{\pi}{2}\)，\(\displaystyle\sum_{n=1}^{\infty}x_n^2\) 收敛，则 \(\displaystyle\prod_{n=1}^{\infty}\cos x_n\) 收敛。

\begin{hintbox}
利用 \(\ln\cos x\sim-\dfrac{x^2}{2}\)（当 \(x\to0\)）。把乘积转成 \(\sum \ln(\cos x_n)\) 再比较。
\end{hintbox}
\end{exercisebox}

\begin{exercisebox}
\textbf{习题 4 \badge{正切乘积的绝对收敛}}

设 \(|a_n|<\dfrac{\pi}{4}\)，\(\displaystyle\sum_{n=1}^{\infty}|a_n|\) 收敛，则 \(\displaystyle\prod_{n=1}^{\infty}\tan\left(\frac{\pi}{4}+a_n\right)\) 绝对收敛。

\begin{hintbox}
可写成 \(\tan(\frac{\pi}{4}+a)=\dfrac{1+\tan a}{1-\tan a}\)，再取对数并利用 \(\tan a\sim a\)。
\end{hintbox}
\end{exercisebox}

\begin{exercisebox}
\textbf{习题 5 \badge{Wallis 型极限}}

证明：

\[
      \begin{aligned}
      &(1)\ \lim_{n\to\infty}\frac{1\cdot3\cdot5\cdots(2n-1)}{2\cdot4\cdot6\cdots(2n)}=0;\\
      &(2)\ \lim_{n\to\infty}\frac{\beta(\beta+1)(\beta+2)\cdots(\beta+n)}
      {\alpha(\alpha+1)(\alpha+2)\cdots(\alpha+n)}=0\qquad(0<\beta<\alpha).
      \end{aligned}
      \]

\begin{hintbox}
第 (1) 直接联想 Wallis 公式。第 (2) 写成无穷乘积或对数和，都能把它压到第 (1) 或积分比较上。
\end{hintbox}
\end{exercisebox}

\begin{exercisebox}
\textbf{习题 6 \badge{无穷乘积恒等式}}

设 \(|q|<1\)，证明：

\[
        \prod_{n=1}^{\infty}(1+q^n)=\frac{1}{\prod_{n=1}^{\infty}(1-q^{2n-1})}.
      \]

\begin{hintbox}
把 \(\prod(1-q^n)\) 分成奇数项和偶数项，再用 \(1-q^{2n}=(1-q^n)(1+q^n)\)。
\end{hintbox}
\end{exercisebox}

\begin{exercisebox}
\textbf{习题 7 \badge{反例构造}}

设

\[
        a_{2n-1}=-\frac1{\sqrt n},
        \qquad
        a_{2n}=\frac1{\sqrt n}+\frac1n\left(1+\frac1{\sqrt n}\right),
      \]

证明级数 \(\displaystyle\sum_{n=1}^{\infty}a_n\) 与 \(\displaystyle\sum_{n=1}^{\infty}a_n^2\) 都发散，但无穷乘积 \(\displaystyle\prod_{n=2}^{\infty}(1+a_n)\) 收敛。

\begin{hintbox}
先把奇偶两项合并，直接算 \((1+a_{2n-1})(1+a_{2n})\)。这题的关键不是分别看 \(a_n\) 和 \(a_n^2\)，而是看成对后的乘积有无望远镜结构。
\end{hintbox}
\end{exercisebox}

来源参考：陈纪修、于崇华、金路《数学分析》第三版下册，第九章 §5。建议把本节和 [第九章 §3 正项级数](./0009-3-positive-series.html)、[第九章 §4 任意项级数](./0009-4-arbitrary-series.html) 连起来看：无穷乘积真正的难点，不在“乘”，而在“如何把它重新送回级数判别法”。

如果你做题后要批改，先发已做部分即可；我会按“必要条件、对数化、判别法、结论”这条线逐题检查。
